\documentclass[10pt,a4paper]{amsart}
\usepackage[margin=19mm]{geometry}
\usepackage{amsmath,amssymb,mathtools}
\usepackage[scr=rsfs,cal=euler]{mathalfa}
\usepackage{xfrac}
\usepackage[unicode,hidelinks,bookmarks=false]{hyperref}
\newcommand{\ie}{i.e.\ }
\newcommand{\cf}{cf.\ }
\newcommand{\resp}{respectively\ }
\newcommand{\Subsection}{\subsection}
\providecommand{\nobreakdash}{\nobreak\hbox{-}}
\setlength{\emergencystretch}{2em}
\multlinegap0pt
\usepackage{bm}
\usepackage{bbm}
\usepackage{dsfont}
\usepackage{xspace}
\usepackage{cancel}
\usepackage{mathtools}
\usepackage{amsthm}
\makeatletter
\@ifundefined{theorem}{\newtheorem{theorem}{Theorem}}{}
\@ifundefined{lemma}{\newtheorem{lemma}[theorem]{Lemma}}{}
\makeatother
\providecommand{\dint}{\displaystyle\int}
\title[On Iterated Lorenz Curves with Applications: The Multivariat]{On Iterated Lorenz Curves with Applications: The Multivariate Case}
\author{Vilimir Yordanov}
\date{Source: arXiv:2507.01406v3 (CC BY 4.0)}
\begin{document}
\maketitle

\noindent\emph{Source and licence.} On Iterated Lorenz Curves with Applications: The Multivariate Case. Version arXiv:2507.01406v3 (CC BY 4.0), distributed under the Creative Commons Attribution 4.0 International licence, as recorded in the arXiv licence metadata for this exact version and shown on the abstract page https://arxiv.org/abs/2507.01406v3. Full rights notice: licenses/arxiv-2507.01406-CC-BY-4.0.txt.\par\smallskip

\noindent\textit{Editorial scope.} Counted unit: the Lemma whose source label is \texttt{Phi} in the article (source0.tex lines 7059--7074, source label \texttt{Phi}), delivered together with the article's own complete original proof of it (source0.tex lines 7076--7261, one \texttt{proof} environment of 186 lines). No mathematics is written, repaired or re-derived: statement and proof are the article's own text line for line. Required context retained uncounted: none. Every definition, display and label of those lines is retained unchanged. Omitted from this package and not redistributed: 1--7058, 7075--7075, 7262--13277. Nothing inside a retained excerpt is omitted or altered: every retained body is delivered exactly as the article's lines are written, so no omission receipt of any kind accompanies this package. Citations: the retained excerpts carry no citation command at all, so no bibliography entry of the article is transcribed and no reference is redistributed. Reference adaptation: none was needed and none was made; every reference inside the delivered unit resolves inside it, so no unresolved reference or citation is produced. Printed slips kept literal: no printed word, symbol, hypothesis or number of the counted statement or of its proof was altered. Layout and class: the article's own class is \texttt{article} with packages including float, capt-of, amssymb, amsmath, geometry, lscape, rotating, setspace, hyperref, cite, appendix, graphicx, caption; the wrapper is the lane's pinned \texttt{amsart} editorial wrapper at 10pt on A4, which restates the theorem-like environments the retained text needs and adds the article's own macro definitions that the retained text needs (\texttt{\textbackslash dint}), with extra wrapper packages (bm, bbm, dsfont, xspace, cancel, mathtools). The delivered numbering is the unit's own. Assets: no retained span contains a figure environment, a graphics-inclusion command, a PDF-page inclusion, a table or a TikZ picture; no image file is copied into this package, the package declares no asset dependency and no image file is redistributed with it. Licence: version arXiv:2507.01406v3 of this article is distributed under the Creative Commons Attribution 4.0 International licence, as recorded in the arXiv licence metadata for this exact version and shown on the abstract page https://arxiv.org/abs/2507.01406v3. A full rights notice is delivered at licenses/arxiv-2507.01406-CC-BY-4.0.txt. Characters: every retained excerpt is pure ASCII, so the delivered engine, XeLaTeX with its default Latin Modern text font, reports no missing character.
\begin{lemma}
\label{Phi} \bigskip For each $n\geq 0$, the function $\Phi _{n}(x)$has two
fixed points at $0$ and $1$, and a unique fixed point in $(0,1)$. If we
denote the latter by $\phi _{n}$ then%
\begin{equation}
\Phi _{n}(x)>x\quad \text{for }x\in (0,\phi _{n}),\qquad \Phi _{n}(x)<x\quad 
\text{for }x\in (\phi _{n},1).  \label{172a1}
\end{equation}

Moreover, there exists a closed interval $I_{n}\subset (0,1)$ such that $%
\phi _{n}\in I_{n}$ and \ \ \ 
\begin{equation}
\Phi _{n}^{^{\prime }}(x)\leq 1\quad \text{for all }x\in I_{n}.
\label{172a2}
\end{equation}
\end{lemma}

\begin{proof}
We can notice that similar behavior to the one described in this claim holds
for the inverse $L_{n}^{-1}(x)$ in terms of the crossing and its general
shape by considering the previous claim and the fact that $L_{n}^{-1}(x)$ is
a mirror image of $L_{n}(x)$ with respect to the diagonal $x$. The current
claim essentially says that the composite function $\Phi _{n}(x)$ inherits
very generally behavior of all the inverses $L_{n}(x)$. It is not at all
obvious how exactly this should happen, and the proof below serves for that
purpose.

We will make the proof by induction. Initially, we consider the composition $%
L^{0,-1}(L^{1,-1}(x))$. We introduce now the deviations $\overline{h}%
_{i}(x)=x-L^{i,-1}(x)$ for $i=0,1$ and the one for the composite function $%
\overline{h}_{\Phi _{1}}(x)=x-\Phi _{1}(x)=x-L^{0,-1}(L^{1,-1}(x)).$ We will
also keep to the notation for the respective fixed points as well as their
properties from the previous claims.

First, we will prove that $\Phi _{1}(x)$ has a unique fixed point in $(0,1)$%
. We notice that we can write%
\begin{equation}
\overline{h}_{\Phi
_{1}}(x)=x-L^{0,-1}(L^{1,-1}(x))=x-L^{1,-1}(x)+L^{1,-1}(x)-L^{0,-1}(L^{1,-1}(x))=%
\overline{h}_{1}(x)+\overline{h}_{0}(L^{1,-1}(x)).  \label{173}
\end{equation}

We can define now $\overline{c}_{0}$ by $\overline{c}_{0}=L_{1}(c_{0})$.
Then for $\overline{h}_{1}$ we have: 
\begin{equation}
\overline{h}_{0}(L^{1,-1}(x))>0\quad \text{for }L^{1,-1}(x)>c_{0}\quad
\Longleftrightarrow \quad \overline{h}_{0}(L^{1,-1}(x))>0\quad \text{for }%
x\in (\overline{c}_{0},1),  \label{174}
\end{equation}%
and 
\begin{equation}
\overline{h}_{0}(L^{1,-1}(x))<0\quad \text{for }L^{1,-1}(x)<c_{0}\quad
\Longleftrightarrow \quad \overline{h}_{0}(L^{1,-1}(x))<0\quad \text{for }%
x\in (0,\overline{c}_{0}).  \label{175}
\end{equation}%
Without loss of generality, we can assume that the unique zero of $\overline{%
h}_{1}$ satisfies $c_{1}<\overline{c}_{0}$, i.e., $c_{1}<L_{1}(c_{0})$ (the
proof for the other case is analogous).

Now we have the following analysis for $\overline{h}_{\Phi _{1}}(x)=%
\overline{h}_{1}(x)+\overline{h}_{0}(L^{1,-1}(x))$ depending on $x:$

\begin{itemize}
\item \textbf{On\ }$(0,c_{1})$\textbf{:} Since $\overline{h}_{1}(x)$ is
negative for $x<c_{1}$ by the crossing property, and since $\overline{h}%
_{0}(L^{1,-1}(x))$ is also negative by the crossing property (because $%
x<c_{1}<\overline{c}_{0}$ implies $L^{1,-1}(x)<L^{1,-1}(c_{1})\leq L^{1,-1}(%
\overline{c}_{0})=c_{0}$), the sum of these two terms must be negative.
Therefore,$\overline{h}_{\Phi _{1}}(x)=\overline{h}_{1}(x)+\overline{h}%
_{0}(L^{1,-1}(x))$ $<0$;

\item \textbf{At }$x=c_{1}$\textbf{:} By definition, $\overline{h}%
_{1}(c_{1})=0$. Since $c_{1}<\overline{c}_{0}$ implies $%
L^{1,-1}(c_{1})<c_{0} $, we have $\overline{h}_{0}(L^{1,-1}(x))<0$. Thus, $%
\overline{h}_{\Phi _{1}}(x)=\overline{h}_{1}(x)+\overline{h}%
_{0}(L^{1,-1}(x)) $ $<0;$

\item \textbf{On }$(c_{1},\overline{c}_{0})$\textbf{:} In this interval, the
two components of $\overline{h}_{\Phi _{1}}(x)$ have opposing signs. Since $%
x>c_{1}$, we have $\overline{h}_{1}(x)>0$; conversely, since $x<\overline{c}%
_{0}$, it follows that $\overline{h}_{0}(L^{1,-1}(x))<0$. The existence of a
root is guaranteed by the \textit{Intermediate Value Theorem}. As
established previously, $h_{\Phi _{1}}(c_{1})<0$. Furthermore, $h_{\Phi
_{1}}(\overline{c}_{0})>0$ (as we will show in the next section). Given this
change in sign, there must be at least one point $\phi _{1}\in $ $(c_{1},%
\overline{c}_{0})$ such that $\overline{h}_{\Phi _{1}}(\phi _{1})=0$. The
uniqueness of this root follows from the strict monotonicity of $\overline{h}%
_{\Phi _{1}}(x)$, a property inherited from $\overline{h}_{0}$ and $%
\overline{h}_{1}$. At this unique point $\phi _{1}$, the positive
contribution from $\overline{h}_{1}$ is precisely balanced by the negative
contribution from $\overline{h}_{0}$;

\item \textbf{At }$x=\overline{c}_{0}$\textbf{:} By definition, $\overline{h}%
_{0}(\overline{c}_{0})=0$. Also, for $x>c_{1}$ we already have $\overline{h}%
_{1}(x)>0$. Thus, $\overline{h}_{\Phi _{1}}(x)=\overline{h}_{1}(x)+\overline{%
h}_{0}(L^{1,-1}(x))$ $>0;$

\item \textbf{On }$(\overline{c}_{0},1)$\textbf{:} For $x>\overline{c}_{0}$,
we have $L^{1,-1}(x)>c_{0}$ and hence $\overline{h}_{0}(L^{1,-1}(x))>0$.
Also, for $x>c_{1}$ we already have $\overline{h}_{1}(x)>0$. Therefore, $%
\overline{h}_{\Phi _{1}}(x)=\overline{h}_{1}(x)+\overline{h}%
_{0}(L^{1,-1}(x))>0.$
\end{itemize}

Second, we will prove that $\Phi _{1}(x)>x$ holds in the interval $(0,\phi
_{1})$ and $\Phi _{1}(x)<x$ holds in the interval $(\phi _{1},1).$ For that
purpose, we want to show that the equation $\overline{h}_{\Phi
_{1}}^{^{\prime }}(x)=0$ has exactly two solutions in $(0,1)$.

Regarding the existence, we already know from above that $\overline{h}_{\Phi
_{1}}(\phi _{1})=0$. This allows to have:

\begin{itemize}
\item On the interval $[0,\phi _{1}]$, since $\overline{h}_{\Phi
_{1}}(0)=0\quad $and$\quad \overline{h}_{\Phi _{1}}(\phi _{1})=0,$ the 
\textit{Rolle's Theorem} guarantees the existence of some $\phi
_{1}^{(1)}\in (0,\phi _{1})$ such that $\overline{h}_{\Phi _{1}}^{^{\prime
}}(\phi _{1}^{(1)})=0;$

\item Similarly, on the interval $[\phi _{1},1]$, since $\overline{h}_{\Phi
_{1}}(\phi _{1})=0\quad $and$\quad \overline{h}_{\Phi _{1}}(1)=0$, there
exists $\phi _{1}^{(2)}\in (\phi _{1},1)$ such that $\overline{h}_{\Phi
_{1}}^{^{\prime }}(\phi _{1}^{(2)})=0.$
\end{itemize}

Regarding the uniqueness, we use a contradiction argument. Take first the
interval $(0,\phi _{1})$. Let's assume that there exist two distinct points $%
a$ and $b$ in it such that $\overline{h}_{\Phi _{1}}^{^{\prime }}(a)=%
\overline{h}_{\Phi _{1}}^{^{\prime }}(b)=0$ holds. Then again by the \textit{%
Rolle's Theorem}, there is a $c\in (a,b)$ with $\overline{h}_{\Phi
_{1}}^{^{\prime \prime }}(c)=0$. But we know that $\overline{h}_{\Phi
_{1}}^{^{\prime }}(x)=\overline{h}_{1}^{^{\prime }}(x)+\overline{h}%
_{0}^{\prime }(L^{1,-1}(x))\,[L^{1,-1}(x)]^{^{\prime }}$ and also $\overline{%
h}_{\Phi _{1}}^{^{\prime \prime }}(x)=\overline{h}_{1}^{^{\prime \prime
}}(x)+\overline{h}_{0}^{\prime \prime }(L^{1,-1}(x))\,\left[
[L^{1,-1}(x)]^{\prime }\right] ^{2}+\overline{h}_{0}^{\prime
}(L^{1,-1}(x))\,[L^{1,-1}(x)]^{^{\prime \prime }}$. Since each deviation
function $\overline{h}_{0}$ and $\overline{h}_{1}$ has exactly one local
maximum on $(0,c_{0})$ and $(0,c_{1})$, their derivatives are strictly
monotone in these intervals. This prevents $\overline{h}_{\Phi
_{1}}^{^{\prime \prime }}(x)$ to be zero and thus we get a contradiction.
Hence, there is at most one solution to $\overline{h}_{\Phi _{1}}^{^{\prime
}}(x)=0$ in $(0,\phi _{1})$. The argument for the interval $(\phi _{1},1)$
is the same.

A direct computation gives%
\begin{eqnarray}
\overline{h}_{\Phi _{1}}^{^{\prime }}(x) &=&1-(L^{0,-1})^{^{\prime
}}(L^{1,-1}(x))(L^{1,-1})^{\prime }(x)=1-\frac{1}{(L_{0}{}^{^{\prime
}})(L^{0,-1}(L^{1,-1}(x)))}\frac{1}{(L_{1}{}^{^{\prime }})(L^{1,-1}(x))} 
\notag \\
&=&1-\frac{\dint\nolimits_{0}^{1}w[F^{-1}(L^{0,-1}(L^{1,-1}(x)))]du}{%
w[F^{-1}(L^{0,-1}(L^{1,-1}(x)))]}\frac{\dint%
\nolimits_{0}^{1}w[L^{0,-1}(L^{1,-1}(x))]du}{w[L^{0,-1}(L^{1,-1}(x))]}.
\label{175.1}
\end{eqnarray}

Hence we have $\overline{h}_{\Phi _{1}}^{^{\prime }}(0)=-\infty $ and $%
\overline{h}_{\Phi _{1}}^{^{\prime }}(1)=-\infty $. We can also notice that
around $0.5$ there is mass concentration which produces at least one
positive value for $h_{\Phi _{1}}^{^{\prime }}(x)$. Therefore at $\phi
_{1}^{(1)}$ and $\phi _{1}^{(2)}$, the function $h_{\Phi _{1}}^{^{\prime
}}(x)$ changes its sign which determine the behavior $\Phi _{n}(x)$ and its
crossing properties. Namely, as hinted at the beginning, the positioning of
the $\Phi _{1}(x)$ above and below the diagonal is analogous to the
situation of $L^{n,-1}(x)$ for any $n$. The crossing pattern is: $\Phi
_{1}(x)>x$ for $x\in (0,\phi _{1})$ and $\Phi _{1}(x)<x$ for $x\in (\phi
_{1},1)\,$, while at each of these intervals the points $\phi _{1}^{(1)}$
and $\phi _{1}^{(2)}$ just produce a maximal distance between $\Phi _{1}(x)$
and the diagonal. Furthermore, the described behavior of $\overline{h}_{\Phi
_{1}}^{^{\prime }}(x)$ gives that $\overline{h}_{\Phi _{1}}^{^{\prime
}}(x)>0,$ and thus $\Phi _{1}^{^{\prime }}(x)<1$, hold for $x\in (\phi
_{1}^{(1)},\phi _{1}^{(2)})$ and $\overline{h}_{\Phi _{1}}^{^{\prime
}}(x)<0, $ and thus $\Phi _{1}^{^{\prime }}(x)>1,$ hold for $x\in \lbrack
0,\phi _{1}^{(1)})\cup (\phi _{1}^{(2)},1]$ with the corresponding zeros
attained at $\phi _{1}^{(1)}$ and $\phi _{1}^{(2)}$. So the interval $I_{1}$
is $[\phi _{1}^{(1)},\phi _{1}^{(2)}]$.

Now we can move to the case $n$. From $\Phi
_{n}(x)=L^{0,-1}(...(L^{n-1,-1}(L^{n,-1}(x)))$ we have $\Phi _{n}(x)=\Phi
_{n-1}(L^{n,-1}(x))$. Either inductively, or simply by change of notation,
we can notice that everything proved above also holds for $\Phi _{n}(x)$.
This is essentially due to the again appearing bifunctional composite form
of $\Phi _{n}(x)$ which was the case also for $\Phi
_{1}(x)=L^{0,-1}(L^{1,-1}(x))$. Exactly, it allows to make the induction
step by assuming for $n$ that $\Phi _{n}(x)$ obeys the conditions of the
claim and then prove them for $n+1$. The proof will verbally reiterate the
proof above for $n=1$ with a change of notation to reflect the fact that
here we work with $\Phi _{n-1}$ instead of $\Phi _{0}\equiv $ $L^{0,-1}$.

We may also finally notice that if we set $\Psi _{n}(x)=\Phi _{n}^{-1}(x)$
then $\Psi _{n}(x)=L_{n}(...(L_{1}(L_{0}(x)))$. The function $\Psi _{n}(x)$
resembles the crossing behavior of $L_{n}(x)$. Namely, $\Psi _{n}(x)$ has
two fixed points at $0$ and $1$. Additionally, it has a unique fixed point
in $(0,1)$. If we denote the latter by $\psi _{n}$, then $\Psi _{n}(x)<x$
holds in the interval $(0,\psi _{n})$ and $\Psi _{n}(x)>x$ holds in the
interval $(\psi _{n},1)$. There is also a closed interval $I_{n}\subset
(0,1) $ such that $\psi _{n}\in I_{n}$ and the derivative satisfies $\Psi
_{n}^{^{\prime }}(x)\geq 1$ for all $x\in I_{n}$.

The proof of the preceding lemma also yields the following immediate
corollary.
\end{proof}

\end{document}
